\section{Love is Blind}

\begin{quotex}
Anyone who does not love does not know God, because God is love.
\flright{\textsc{1 John 4:8}}

The proof of love is the showing forth of the work.
\flright{\textsc{Gregory the Great}, \emph{Homily XXX}}
\end{quotex}


You may have heard it said that the Intellect is superior to Love, because the Intellect knows, but Love is blind. The opposite is the case. The Intellect is passive because it can only reflect the Truth, but Love is Active because it acts and can create the Good.

Thus Plato stopped short at the contemplation of the Idea of the Good, mistaking that for God. Since God is Love, the next stage is to bring the Good into manifestation, as Gregory the Great said. Nevertheless, Plato was correct because one should only love what is good, and only the Intellect, or nous, knows the Good. Therefore, the role of the Intellect is to show Love what is worthy to be loved.

\begin{itemize}


\item Will of Heaven: Upward Pull

\item Mystical marriage

\item Nous: Free choice of the Good

\item Thumos: Emotional Attachment

\item Eros: Blind Attraction

\item Physical: Fundamental forces

\item Gravity: Downward pull to matter
\end{itemize}


Love is manifested differently in each of the states of being. At the physical layer, attraction is automatic and indiscriminate. Magnets attract just because they are magnetic. At the lower levels of the human being, eros and thumos, Love is truly blind. Eros drives the sexual instinct without regard to the goodness of the object of desire. Thumos adds a strong emotional content to desire. Emotions make Love blind to the idea of goodness or badness. We see that frequently in relationships. Even after the loved one is bad, emotions make it difficult the end things.

Nous, on the other hand, is objective and can distinguish between the good and bad. In an ordered soul, the knowledge of the nous will dominate thumos and eros: it will desire what is good and feel a corresponding attraction. Otherwise, desires and emotions will overwhelm the Intellect so that it serves man's lower nature instead of leading the being to seek the Good.

\paragraph{Alchemical Marriage}


\begin{quotex}
For whoever has achieved the perfect realization of unity in himself, all opposition has ceased and with it the state of war, for from the standpoint of totality, which lies beyond all particular standpoints, nothing remains but absolute order. Nothing can thereafter harm such a being, since for him there are no longer any enemies, either within him or without.
\flright{\textsc{Rene Guenon}, \emph{Symbolism of the Cross}}

The psychic preparation described by Dorn an alchemist is therefore an attempt, uninfluenced by the East, to bring about a union of opposites in accordance with the great Eastern philosophies, and to establish for this purpose a principle freed from the opposites and similar to the atman or tao.
\flright{\textsc{Carl Jung}, \emph{Mysterium Coniunctionis}}
\end{quotex}


The highest good of the human state is to achieve wholeness. The alchemists described a path to that realization, albeit cloaked in symbolic language. It begins with self-knowledge. It is imperative to emphasize that knowledge of the Self is NOT the same as knowledge of the ego or one's personality quirks.

That wholeness is a return to the Primordial State, the union of Adam-Eve. The metals of alchemical symbolism refer to the male and female aspects of the soul. For a man, the love of a woman, real or imagined, who transcends the duties of human marriage, helps awaken that self-knowledge. Ultimately, they will unite in a higher state of being.


\flrightit{Posted on 2022-07-31 by Cologero}

\begin{center}* * *\end{center}

\begin{footnotesize}
\begin{sffamily}
\texttt{Ole on 2022-08-01 at 07:40 said:}

The Wedding at Cana :-)

\hfill

\texttt{Adil on 2022-08-01 at 12:34 said:}

One has to wonder whether the metaphysical cause to the gender dysphoria of the current age is a result of the disintegration of this Adam-Eve principle, where the poles come in conflict. That would mean that there is an ironic merit to the gender fluid aspirations of postmodernism, albeit not in the sense of physically changing your gender, but actually achieving the ``non-binary'' wholeness that comes from integrating the anima or animus.

The feminine heart of the Catholic Church once greatly attracted me to her, which seemed to have a mystical lure that simply was gone in the Protestant Christianity. I realize now it was the Virgin, who has been neglected in the West.

\hfill

\texttt{Auld Wat on 2022-08-01 at 17:16 said:}

Will of Heaven: Upward Pull

Mystical marriage

Nous: Free choice of the Good

Thumos: Emotional Attachment

Eros: Blind Attraction

Physical: Fundamental forces

Gravity: Downward pull to matter

Is the first part missing from no. 2?

\hfill

\texttt{Balder on 2022-08-02 at 17:55 said:}

@Adil: Robert Bolton has a very good phrase in his The Order of the Ages that can be related to tthe gender dysphoria: The Submergence of Distinctions. In this he describes how Man and Woman are archetypes incarnated in human beings, and I would say that the gender dysphoria some increasingly feel is one sign of the times in which the divine archetypes are increasingly under dissolution from inside and outside. I regard the propaganda related to the issue to be almost absolutely evil in which a mental illness related to the relentless dissolution is given an almost holy appearance, as in the sense of ``cursed saints of our times''. I can easily recommend the book as a whole to see what is going on and how universal time realizes evil. The ``submergence of distinctions'' is a one way to say how the diabolical impulse can realize the perversion of divine unity.


\hfill

\end{sffamily}
\end{footnotesize}

